% ---------------------------------------------------------------------
%  TÓM TẮT
% ---------------------------------------------------------------------
\begin{abstract}

Tìm phòng trọ phù hợp là nhu cầu thiết yếu của sinh viên và người lao động
tại Đà Nẵng, nhưng thông tin cho thuê hiện nay nằm rải rác trên nhiều kênh,
thiếu công cụ lọc theo đúng nhu cầu cá nhân về ngân sách, khu vực và tiện
ích. Đồ án xây dựng một hệ thống hỗ trợ tìm phòng trọ tại Đà Nẵng, trong đó
thành phần cốt lõi là một hệ gợi ý theo hướng dựa trên nội dung
\en{(content-based filtering)}, kết hợp thêm tín hiệu từ những người dùng có
nhu cầu tương tự để tăng độ đa dạng của gợi ý. Hệ thống trích xuất đặc trưng
của từng phòng trọ gồm giá thuê, khu vực, diện tích, tiện ích và mô tả văn
bản, sau đó tính độ tương đồng có trọng số với hồ sơ nhu cầu người dùng để
xếp hạng kết quả. Dữ liệu thử nghiệm gồm 2.500 tin đăng phòng trọ thu thập
tại sáu quận trung tâm của Đà Nẵng, cùng lịch sử tương tác của 200 người
dùng thử nghiệm dùng để đánh giá ngoại tuyến. Hệ thống đạt Precision@5 bằng
0,76 và NDCG@5 bằng 0,81, cao hơn mô hình chỉ dùng nội dung thuần tuý lần
lượt 7 và 6 điểm phần trăm. Khảo sát nhanh trên 30 sinh viên dùng thử cho
mức độ hài lòng trung bình 4,1 trên thang 5. Toàn bộ hệ thống được triển
khai thành ứng dụng web cho phép người dùng nhập tiêu chí và nhận danh sách
phòng trọ gợi ý kèm bản đồ vị trí. Kết quả cho thấy cách tiếp cận dựa trên
nội dung là lựa chọn khả thi khi dữ liệu tương tác người dùng còn hạn chế,
và có thể mở rộng thêm thành phần cộng tác khi hệ thống có nhiều người dùng
hơn.

\vspace{0.8em}
\noindent\textbf{Từ khoá:} hệ gợi ý, lọc dựa trên nội dung, độ tương đồng
cô-sin, phòng trọ, Đà Nẵng.

\end{abstract}
